% !TeX spellcheck = ru_RU
% !TEX root = vkr.tex

\section*{Введение}
\thispagestyle{withCompileDate}

Реляционная модель данных, предложенная Коддом~\cite{codd1970relational}, остаётся доминирующей парадигмой организации информации. Несмотря на распространение альтернативных подходов~\cite{davoudian2018survey}, реляционные СУБД сохраняют ведущие позиции в индустрии благодаря зрелой экосистеме и развитым механизмам обеспечения целостности~\cite{abiteboul1995foundations}.

Однако с ростом объёмов данных схемы баз данных неизбежно деградируют: ограничения устаревают, а реальные закономерности перестают соответствовать первоначально заданным правилам~\cite{abedjan2015profiling}. В связи с этим возрастает потребность в профилировании данных~--- систематическом анализе содержимого баз данных для извлечения метаданных~\cite{naumann2014data, abedjan2018data}. Результаты профилирования применяются при оптимизации запросов~\cite{chaudhuri1998overview}, обнаружении аномалий и повышении качества данных~\cite{rahm2000data} по некоторым заранее выбранным метрикам.

Центральным инструментом профилирования служат зависимости данных~--- формально описанные закономерности в хранимой информации~\cite{liu2012discover}. Среди них важное место занимают зависимости включения (Inclusion Dependencies, IND)~\cite{casanova1984inclusion, de2002unary}: отношение между атрибутами таблиц, при котором множество значений одного атрибута является подмножеством значений другого. Зависимости включения играют ключевую роль при обнаружении внешних ключей~\cite{rostin2009framework} и поддержании ссылочной целостности~\cite{levene2000justification}.

Автоматический поиск зависимостей включения является NP-трудной задачей~\cite{kantola1992discovering}, поэтому для её решения был разработан ряд специализированных алгоритмов: SPIDER~\cite{bauckmann2006efficiently}, BINDER~\cite{papenbrock2015divide}, MIND~\cite{de2002unary} и FAIDA~\cite{kruse2018fast}. Тем не менее на практике зависимости включения далеко не всегда выполняются на всех данных. Для таких случаев были введены приближённые зависимости включения (Approximate IND, AIND)~\cite{kruse2018fast}, допускающие определённую долю нарушений. Однако приближённые зависимости дают лишь агрегированную оценку доли нарушений, не раскрывая того, на каком именно подмножестве данных зависимость выполняется.

Эту проблему решают условные зависимости включения (Conditional Inclusion Dependencies, CIND), предложенные Браво и Фанем~\cite{bravo2007extending}. CIND расширяют классические IND добавлением таблицы условий~--- набора шаблонов, определяющих подмножество кортежей, на котором зависимость должна выполняться. Такой подход позволяет выявлять нарушения ссылочной целостности на определённых срезах данных~\cite{fan2008conditional, brink2012cind}, оптимизировать выполнение запросов~\cite{golab2008generating}, а также обнаруживать закономерности, проявляющиеся только при определённых условиях~\cite{chiang2008discovering}.

Для поиска условий CIND в~\cite{brink2012cind} были предложены два алгоритма: CINDERELLA, основанный на ассоциативных правилах с использованием адаптированного алгоритма Apriori~\cite{agrawal1994fast}, и PLI-CIND, работающий с пересечением позиционных списков атрибутов. CINDERELLA осуществляет поиск кандидатов в ширину, обеспечивая высокую скорость ценой повышенного потребления памяти, тогда как PLI-CIND использует поиск в глубину с существенно меньшими требованиями к памяти.

Алгоритмы профилирования данных реализованы в нескольких инструментах, наиболее известными из которых являются Metanome~\cite{papenbrock2015data} и Desbordante~\cite{smirnov2025desbordante}. Desbordante~--- высокопроизводительный инструмент с открытым исходным кодом, написанный на C++ и предоставляющий интерфейсы через Python и командную строку.

Ранее в рамках бакалаврской работы~\cite{abrosimov2025cinderella} была выполнена первоначальная реализация алгоритмов CINDERELLA и PLI-CIND и их интеграция в Desbordante. Однако эта реализация имеет ряд ограничений: алгоритмы не были оптимизированы по производительности, а задача верификации~--- проверки того, выполняется ли заданная CIND на конкретных данных,~--- не рассматривалась. Между тем верификация является важной практической задачей, так как позволяет проверять гипотезы о закономерностях без запуска полного поиска~\cite{papenbrock2016functional, bleifuss2024efficient}.

Настоящая работа посвящена оптимизации существующих алгоритмов поиска условных зависимостей включения, а также разработке и реализации алгоритма верификации CIND в инструменте профилирования данных Desbordante.

\section{Постановка задачи}
\label{sec:task}
Целью работы является оптимизация алгоритмов поиска условных зависимостей включения, разработка алгоритма верификации CIND и их интеграция в инструмент профилирования данных Desbordante с последующим экспериментальным исследованием.

Для достижения поставленной цели необходимо решить следующие задачи:
\begin{enumerate}
    \item Провести оптимизацию существующих реализаций алгоритмов поиска условий CINDERELLA и PLI-CIND, подготовить набор тестов и примеры использования.
    \item Разработать и реализовать алгоритм верификации CIND, позволяющий проверять выполнимость заданной условной зависимости включения на конкретных данных, интегрировать его в Desbordante с тестами и примерами использования.
    \item Провести экспериментальное исследование: подтвердить корректность оптимизированных алгоритмов, оценить влияние оптимизаций на время работы и потребление памяти, а также измерить производительность верификатора CIND.
\end{enumerate}
